% Chapter 7: Refactoring Legacy Projects with Quality Control
\chapter{Refactoring Legacy Projects with Quality Control}
\label{ch:refactoring}

Introducing Claude to an existing codebase is harder than starting fresh.
The codebase has conventions Claude must learn, technical debt it must
navigate, and invariants it must not break. This chapter covers the
protocol for safe, incremental improvement.

\section{Brownfield: Introducing Claude to an Existing Codebase}
\label{sec:brownfield}

\marginnote[Official]{Use Plan Mode first---let Claude read and understand the codebase without modifying anything.\sourceurl{https://code.claude.com/docs/en/common-workflows}}

\subsection{Start with Understanding, Not Action}

\official{Before any changes, use Plan Mode (Shift+Tab) to let Claude explore.}
Ask:
\begin{itemize}
  \item ``Explain the architecture of this project.''
  \item ``What patterns does this codebase use?''
  \item ``Where are the tests? What is the test coverage?''
  \item ``What are the main entry points?''
\end{itemize}

\official{Document existing conventions in CLAUDE.md, not aspirational ones.}
A CLAUDE.md that says ``we use strict TypeScript'' when the codebase has
\code{any} everywhere will cause Claude to produce inconsistent code---
some files strict, some loose.

\subsection{Incremental Adoption}

\practitioner{Adopt Claude configuration in four incremental steps,
each independently valuable:}

\begin{enumerate}
  \item \textbf{CLAUDE.md} documenting current state (build commands, conventions as they are)
  \item \textbf{Permissions} restricting Claude to safe operations
  \item \textbf{Hooks} enforcing existing standards (not new ones)
  \item \textbf{Skills} emerging organically as patterns stabilize
\end{enumerate}

% -------------------------------------------------------------------
\section{The Incremental Refactoring Protocol}
\label{sec:refactoring-protocol}

\practitioner{Never attempt a big-bang rewrite.}
Use four-step increments, each independently shippable:

\begin{description}
  \item[1. Extract] Isolate the function or module to refactor.
    Create a clear boundary between ``changing'' and ``not changing.''

  \item[2. Test] Write characterization tests capturing current behavior.
    Claude excels at this---prompt: ``Write tests that document exactly
    what this function does today, including edge cases and error handling.''

  \item[3. Harden] Refactor with tests as safety net.
    Each change verified against the characterization tests.
    If a test fails, you have broken existing behavior.

  \item[4. Promote] Move refactored code into production.
    Run full regression suite. Verify no downstream breakage.
\end{description}

\begin{tipbox}[Session Sizing]
Each step fits in a 25-minute focused session.
Commit after each step.
If interrupted, the project is always in a working state.
\end{tipbox}

\subsection{Writer/Reviewer Pattern}

Use separate Claude sessions for writing and reviewing
to avoid confirmation bias.
The writer session has context about the refactoring intent and may
overlook issues. The reviewer session sees only the diff and judges
it on its own merits.
See \cref{ch:parallel} for the full treatment of this pattern.

% -------------------------------------------------------------------
\section{Quality Control During Refactoring}
\label{sec:refactoring-quality}

\subsection{Characterization Tests First}

\convergence{Before changing any code, have Claude write tests that capture
existing behavior.} These become your regression safety net.

\begin{minted}{text}
Prompt: "Read src/payments/legacy_processor.py.
Write pytest tests that document its exact current behavior:
- What inputs does it accept?
- What does it return for each input type?
- What errors does it raise and when?
- What side effects does it produce?
Do NOT write tests for how it SHOULD work.
Write tests for how it DOES work."
\end{minted}

\margintip{Characterization tests are temporary---they may test buggy behavior. That is intentional. They exist to detect \emph{changes}, not to assert \emph{correctness}.}

\subsection{Metrics-Driven Approach}

\practitioner{Establish baseline metrics before refactoring.
Track improvement per sprint. Make progress visible with traffic-light thresholds.}

\begin{fullwidth}
\begin{tabular}{@{}llll@{}}
\toprule
\textbf{Metric} & \textbf{GREEN} & \textbf{YELLOW} & \textbf{RED} \\
\midrule
Test coverage & $\geq$80\% & 60--79\% & $<$60\% \\
Lint warnings & 0 & 1--10 & $>$10 \\
Avg function length & $\leq$30 lines & 31--50 & $>$50 \\
Presentation pass rate & $\geq$95\% & 85--94\% & $<$85\% \\
\bottomrule
\end{tabular}
\end{fullwidth}

\subsection{Guard Rails via Hooks}

\practitioner{Hooks protect the refactoring process:}

\begin{description}
  \item[PreCommit] Run full test suite. Block commits if tests fail.
  \item[PostFileWrite] Auto-format to maintain style consistency.
  \item[Custom skill] Run coverage diff: ``Did this change reduce coverage?''
\end{description}

\subsection{Impact Analysis with Subagents}

\official{Use subagents for impact analysis:
``What other files depend on this module?''
runs in isolated context without polluting your main session.}

% -------------------------------------------------------------------
\section{Migration Patterns at Scale}
\label{sec:migrations}

For bulk migrations across many files, use the Fan-Out pattern
described in \cref{ch:automation}. For parallel refactoring across
independent modules, use git worktrees as described in \cref{ch:parallel}.

\subsection{The Quality Flywheel}

\practitioner{After each refactoring sprint:}

\begin{enumerate}
  \item Update CLAUDE.md with lessons learned
  \item Add new hooks for discovered gotchas
  \item Refine skills based on what worked
  \item Update metrics baselines
\end{enumerate}

Each sprint improves not just the code but the tooling around the code.
After three sprints, the refactoring process itself is significantly faster.

% -------------------------------------------------------------------
\section{Visual: The Incremental Refactoring Cycle}

\begin{figure}[H]
\centering
\begin{tikzpicture}[
  step/.style={draw=PrimaryBlue, fill=PrimaryBlue!8, rounded corners=6pt,
    minimum width=2.5cm, minimum height=1.2cm, font=\small\bfseries, align=center},
  note/.style={font=\scriptsize, text=DarkText, align=center},
  arr/.style={->, very thick, PrimaryBlue!60},
  >=Stealth
]
  % Four steps in a cycle
  \node[step] (extract) at (0,2.5) {1. Extract\\[-2pt]\normalfont\scriptsize Isolate target};
  \node[step] (test) at (3.5,2.5) {2. Test\\[-2pt]\normalfont\scriptsize Characterize};
  \node[step] (harden) at (3.5,0) {3. Harden\\[-2pt]\normalfont\scriptsize Refactor};
  \node[step] (promote) at (0,0) {4. Promote\\[-2pt]\normalfont\scriptsize Ship};

  % Arrows forming cycle
  \draw[arr] (extract) -- (test);
  \draw[arr] (test) -- (harden);
  \draw[arr] (harden) -- (promote);
  \draw[arr] (promote) -- (extract) node[midway, left, note] {Next\\module};

  % Center label
  \node[font=\footnotesize\itshape, text=TipGreen!80!black, align=center] at (1.75,1.25)
    {Commit after\\each step};

  % Time annotation
  \node[font=\scriptsize, text=PrimaryBlue!70] at (1.75,-0.8)
    {Each step $\approx$ 25 min $\cdot$ Always shippable};
\end{tikzpicture}
\caption{The incremental refactoring cycle. Each step is independently shippable.}
\label{fig:refactoring-cycle}
\end{figure}

\begin{keyinsight}
Refactoring with Claude follows the same principle as refactoring without it:
small, verifiable increments. The difference is that Claude can write
characterization tests in minutes instead of hours, making the
Extract $\rightarrow$ Test $\rightarrow$ Harden $\rightarrow$ Promote cycle
dramatically faster.
\end{keyinsight}
